Na configuração com apenas cinco sensores MQ, o \textit{Extra Trees} errou 2.932 das 52.486 leituras (5,59\%). Desses erros, 2.700 ocorreram em vasos inoculados e 232 em controles. Cinco coletas de vasos inoculados concentraram 92,09\% dos erros, sobretudo as coletas do vaso 8 no dia 19 (754/1.317 leituras), do vaso 7 no dia 11 (627/2.048) e do vaso 9 no dia 18 (544/1.327). Os erros formaram trechos temporais dentro dessas coletas: no vaso 7 do dia 11, 96,5\% das leituras do último quarto da coleta foram classificadas incorretamente; no vaso 9 do dia 18, todas as leituras desse quarto foram classificadas incorretamente. Entre os erros, 1.315 (44,85\%) tiveram probabilidade prevista para a classe errada entre 0,50 e 0,60, enquanto 166 (5,66\%) tiveram probabilidade de pelo menos 0,90; estes últimos ocorreram todos no trecho final da coleta do vaso 7 no dia 11. As probabilidades não foram calibradas e, portanto, indicam a distância do limiar de decisão, não a certeza estatística do classificador. Após a média das probabilidades por coleta, os cinco MQ classificaram corretamente as 31 coletas e os dez vasos, mas a coleta do vaso inoculado 8 no dia 19 ficou próxima do limiar (probabilidade média de 0,4984 para a classe 1). Esses padrões são exploratórios e podem orientar o controle do tempo de coleta, das condições ambientais e da deriva dos sensores em experimentos com novos vasos.
